\section{Modelamiento del Problema}

\subsection{Formulación del problema de detección}
\label{sec:formalizacion}

Sea $r_i(t) = \ln P_i(t) - \ln P_i(t-1)$ el retorno logarítmico del token $i$
en el período $[t-1, t]$. Sobre una ventana de tiempo de longitud $T$,
representamos la dinámica conjunta del mercado mediante la matriz de correlación
$C_t \in \mathbb{R}^{N \times N}$ con entradas:
\begin{equation}
  C_{ij}(t) = \frac{\langle r_i r_j \rangle - \langle r_i \rangle\langle r_j \rangle}
  {\sigma_i \,\sigma_j},
  \label{eq:pearson}
\end{equation}
donde $\sigma_i$ es la desviación estándar de los retornos del token $i$. Esta
matriz captura el grado de dependencia estadística entre pares de tokens en
cada ventana temporal.

Modelamos un evento de \emph{pump \& dump} como la aparición transitoria de un
subconjunto de tokens $S \subset V$ (donde $V$ es el conjunto de nodos del
grafo de correlaciones, definido formalmente en la Sección~\ref{sec:marco}),
con $|S| \ll N$, cuyos retornos presentan
dependencia estadística anómala respecto al resto del mercado. La acumulación
coordinada previa a un \emph{pump} induce correlaciones elevadas dentro de $S$,
generando una estructura colectiva que no está presente en condiciones normales.

Bajo esta perspectiva, el problema se formula como una tarea de detección:
identificar ventanas temporales $t$ en las que existe un subconjunto $S$ tal
que las correlaciones entre sus elementos son sistemáticamente mayores ---o
estructuralmente distintas--- que las esperadas bajo un comportamiento no
manipulado. Este enfoque no requiere conocer a priori qué tokens pertenecen a
$S$, sino inferir su existencia a partir de patrones agregados en $C_t$.

\subsection{Modelo formal}
\label{sec:modelo}

A partir de la matriz de correlación $C_t$ definida en la
ecuación~(\ref{eq:pearson}), el pipeline de filtrado opera sobre el espectro de
$C_t$ para separar señal de ruido y construir la red de correlaciones.

El problema es que esta matriz está llena de ruido: la mayoría de las
correlaciones no reflejan una relación real entre tokens, sino fluctuaciones
aleatorias. Para separar señal de ruido aplicamos la Teoría de Matrices
Aleatorias (RMT). Descomponemos $C_t = \sum_{k=1}^{N} \lambda_k
\mathbf{v}_k \mathbf{v}_k^\top$ y usamos la distribución de
Marchenko-Pastur~\cite{plerou1999}, que predice el intervalo de eigenvalores
del ruido puro:
\begin{equation}
  \lambda_{\pm} = \sigma^2\!\left(1 \pm \sqrt{Q}\,\right)^{\!2},
  \label{eq:mp}
\end{equation}
donde $Q = N/T$ y $\sigma^2$ es la varianza promedio de los retornos.
Los eigenvalores en $[\lambda_-, \lambda_+]$ son ruido y se descartan.
El eigenvalor dominante $\lambda_1$, que corresponde al modo de mercado
global, también se elimina.

La RMT requiere $T \gg N$ para ser válida. Un \emph{pump \& dump} dura
solo horas~\cite{lamorgia2020}, lo que obliga a usar ventanas de detección
cortas donde esa condición no se cumple. La solución es separar los dos
pasos: estimar los límites con una ventana larga de calibración
$T_{\text{cal}} = 720~\text{h}$ (30~días) y aplicarlos sobre ventanas
cortas de $T_{\text{det}} = 48~\text{h}$. La validez de la transferencia
se verifica mediante bootstrap: se extraen $B = 1000$ ventanas de
calibración no solapadas del histórico disponible, se calcula
$\lambda_+^{(b)}$ en cada una, y se comprueba que la precisión del
detector varía menos del $5\%$ entre los percentiles~5 y~95 de esa
distribución. Este procedimiento captura la variabilidad real de
$\lambda_+$ bajo los distintos regímenes de volatilidad observados en el
mercado, sin asumir un rango de perturbación arbitrario.

La matriz filtrada que nos queda es:
\begin{equation}
  \tilde{C}_{ij} = \sum_{k \in \mathcal{S}} \lambda_k\, v_{ki}\, v_{kj},
  \label{eq:Ctilde}
\end{equation}
donde $\mathcal{S} = \{k : \lambda_k > \lambda_+,\; k \neq 1\}$ contiene solo
los eigenvalores informativos, sin el ruido ni el modo de mercado.

Con esta matriz construimos la red. El peso de cada arista $(i,j)$ se obtiene
de la distancia de Mantegna~\cite{mantegna1999}:
\begin{equation}
  d_{ij}(t) = \sqrt{2\left(1 - \tilde{C}_{ij}(t)\right)},
  \label{eq:mantegna}
\end{equation}
con $w_{ij} = 2 - d_{ij}$: tokens más correlacionados tienen aristas más
pesadas. Esta métrica tiene la ventaja de preservar tanto correlaciones
positivas como negativas, a diferencia de simplemente truncar los valores
negativos a cero.

Si usáramos todos los pares de tokens como aristas, la red sería
demasiado densa: cada nodo conectado con todos los demás hace que los
eigenvectores de la Laplaciana sean uniformes y el IPR quede cercano a
$N^{-1}$ para todos los modos, borrando cualquier señal de localización. Para
evitar esto, reducimos la red en dos pasos: construimos el Árbol de
Expansión Mínima (MST) sobre las distancias $d_{ij}$, que conecta todos
los nodos con $N-1$ aristas preservando las correlaciones más fuertes;
luego añadimos las $k$ aristas de mayor peso faltantes en cada nodo,
con $k$ elegido como el valor mínimo tal que la similitud del coseno
promedio entre los primeros 20 eigenvectores del grafo completo y del
grafo esparsificado supere $\bar{s}_\phi \geq 0.95$.

La conservación de la estructura espectral se verifica en dos niveles:
(i) los primeros 20 eigenvalores antes y después de la reducción tienen
correlación de Spearman $> 0.9$; (ii) $\bar{s}_\phi > 0.95$ por
construcción. El segundo criterio es esencial porque el IPR depende de
$\phi_{ki}^4$ --- por el teorema de Davis-Kahan, eigenvectores
correspondientes a eigenvalores casi degenerados pueden rotar
significativamente ante una perturbación aun cuando los eigenvalores
apenas cambien. Verificar solo eigenvalores no garantiza la estabilidad
del estadístico de detección.

Sobre el grafo reducido $G_t = (V, E, w)$ calculamos la Laplaciana
normalizada:
\begin{equation}
  \mathcal{L}_t = D_t^{-1/2}(D_t - A_t)\,D_t^{-1/2},
  \label{eq:laplacian}
\end{equation}
donde $A_t$ es la matriz de adyacencia y $D_t$ la matriz diagonal de grados.
Sus pares propios $\{\lambda_k, \boldsymbol{\phi}_k\}_{k=0}^{N-1}$ satisfacen
$0 = \lambda_0 \leq \lambda_1 \leq \cdots \leq \lambda_{N-1} \leq 2$.
El eigenvalor $\lambda_0 = 0$ corresponde al vector constante y se excluye
del análisis.

Para cada eigenvector $k \geq 1$ calculamos el \emph{Inverse Participation
Ratio}:
\begin{equation}
  \mathrm{IPR}_k(t) = \sum_{i=1}^{N} \phi_{ki}^4,
  \label{eq:ipr}
\end{equation}
donde $\phi_{ki}$ es la componente $i$ del eigenvector $k$. Un valor
$\mathrm{IPR}_k \approx N^{-1}$ indica participación uniforme de todos
los tokens (eigenvector extendido); un valor cercano a~$1$ indica
concentración en un solo nodo. Para un grupo artificial de $m$ tokens,
el IPR escala como ${\sim}\,m^{-1}$.

La posición espectral donde aparece la localización de un grupo manipulado
depende de cuántos tokens lo componen y cómo se conectan al resto de la
red. Sin un modelo previo que lo indique, restringir el análisis a los
primeros o últimos eigenvalores no está justificado.

\subsubsection*{El problema del \emph{eigenvalue crossing}}

Una forma ingenua sería seguir el IPR del eigenvector número $k$ a lo largo
del tiempo. El problema es que en redes dinámicas el orden de los eigenvalores
cambia entre ventanas consecutivas: el eigenvector número 5 hoy puede
corresponder a una estructura completamente distinta que el número 5 de mañana.
Comparar $\mathrm{IPR}_5(t)$ con $\mathrm{IPR}_5(t+1)$ puede ser como comparar
cosas distintas, produciendo una serie temporal sin sentido.

\subsubsection*{Solución: resumir todo el espectro}

En lugar de seguir modos individuales, en cada ventana calculamos dos
estadísticos sobre la distribución completa de IPR:
\begin{align}
  \mathrm{IPR}_{\text{mean}}(t) &= \frac{1}{N-1}
    \sum_{k=1}^{N-1} \mathrm{IPR}_k(t), \label{eq:iprmean} \\[3pt]
  \mathrm{IPR}_{p90}(t) &= \mathrm{percentil}_{90}
    \bigl\{\mathrm{IPR}_k(t) : k = 1, \ldots, N-1\bigr\}. \label{eq:iprp90}
\end{align}
Este estadístico no depende del orden de los eigenvalores, así que es
inmune al problema del \emph{eigenvalue crossing}. La elección del
percentil~90 balancea sensibilidad y robustez: el promedio
$\mathrm{IPR}_{\text{mean}}$ diluye la señal cuando la manipulación
afecta a $m \ll N$ tokens --- para $m = 5$ y $N = 50$, el modo
localizado ($\mathrm{IPR} \sim 1/5$) promediado con los $N-2$ modos
extendidos ($\mathrm{IPR} \sim 1/N$) produce un promedio apenas un $4\%$
mayor que el baseline, indetectable; el $\mathrm{IPR}_\text{max}$ o
percentiles extremos (p99) son más ruidosos ante modos
esporádicamente localizados no relacionados con manipulación. Para un
grupo de $m \geq 5$ tokens, el modo localizado cae sistemáticamente en
el decil superior del espectro de IPR. La sensibilidad a esta elección
se evalúa en el análisis de robustez comparando p80, p90, p95 y p99.

La hipótesis es la siguiente: cuando un grupo de tokens está siendo acumulado
coordinadamente antes de un \emph{pump}, forman una comunidad artificial en la
red que hace que al menos un modo espectral se concentre en ellos. No necesitamos
saber cuál modo es ese, solo detectar que existe un aumento anómalo en
$\mathrm{IPR}_{p90}(t)$.

Con este marco, la pregunta central planteada al inicio se formaliza como:
\begin{quote}
  \emph{¿Puede el percentil~90 de la distribución del IPR de la Laplaciana
  normalizada de una red de correlaciones filtrada detectar la fase de
  acumulación previa a esquemas de \emph{pump \& dump} en criptoactivos?}
\end{quote}

Si ese aumento no aparece en los eventos confirmados, el resultado negativo
también es informativo y publicable.
